\subsection{Properties of Power Functions} 
A \emph{power function} is a function in the form
\[
f(x)=x^p
\]
for some exponent $p$. We also allow them to be composed with \makeblank[1in] functions, so the following functions would also be considered power functions:
\[
f(x)=2x^{5/3}-1
\qquad
f(x)=-3(x+4)^7-8
\qquad
f(x)=(2x-5)^{1/2}+4
\]
The properties of power functions depend on the exponent. If the exponent is rational, write $p=\frac{m}{n}$ in lowest terms. In this case, we can rewrite 
$$f(x)=x^{m/n}=\makeblank$$
\begin{itemize}
\item If $p$ is irrational, $x^p$ is defined only for \makeblank[1.5in] numbers.
\item If $p$ is negative, $x^p$ is undefined at \makeblanksmall, and it's never equal to \makeblanksmall.
\item If $p$ is rational and $m$ is even, then $f(x)$ is always \makeblank[1.5in], since we are taking an even \makeblank[1in].
\item If $p$ is rational and $n$ is even, then $f(x)$ is defined only for \makeblank[1.5in] numbers, since we are taking an even \makeblank[1in]. We also know that even \makeblank[1in] are always \makeblank[1.5in]. 
\item With these exceptions, $x^p$ has range and domain of all real numbers.
\end{itemize}
\begin{example}
Find the range and domain of $f(x)=x^{-5/2}$
\begin{itemize}
\item $p=\makeblanksmall$; this is a \makeblank[1.5in] number.
\item $p$ is \makeblank[1.5in], so \makeblanksmall is not in the domain or range.
\item $n$ is \makeblank[1in], so only \makeblank[1.5in] numbers are in the domain.
\item Also because $n$ is \makeblank[1in], only \makeblank[1.5in] numbers are in the range.
\item These are all the restrictions, so:
\end{itemize}
\[
\dom f = \makeblank[1in]
\qquad
\rng f = \makeblank[1in]
\]
\end{example}